problem:  
Let \((M^n,g,f)\) be a **complete noncompact gradient shrinking Ricci soliton** satisfying  
\[
\operatorname{Ric} + \nabla^2 f = \tfrac12 g,
\]
and suppose that \(M\) is **κ–noncollapsed on all scales** and has **nonnegative sectional curvature**. Assume moreover that \(M\) is **not** a Riemannian product \(N^k\times\mathbb{R}^{n-k}\) with a compact Einstein factor \(N^k\).  
Define for each \(r>0\) the rescaled manifold  
\[
(M, r^{-2} g, p),
\]
where \(p\) is a minimum point of \(f\). When \(r_i\to\infty\), these rescalings have, after passing to a subsequence, pointed Gromov–Hausdorff limits \((M_\infty,g_\infty,p_\infty)\).  

It is known that for several special shrinkers (such as the rotationally symmetric ones), the limit \(M_\infty\) is a **unique metric cone with smooth Einstein cross‑section**. However, general uniqueness of the asymptotic cone is not known, even under nonnegative curvature.  

**Question / Conjecture:**  
For a nonflat, noncollapsed gradient shrinking Ricci soliton with nonnegative sectional curvature and no Euclidean factor,  

1. Is the blow‑down limit \((M_\infty,g_\infty)\) **unique** (independent of the chosen subsequence \(r_i\to\infty\))?  
2. If not, can one quantitatively relate all possible blow‑down limits by a **canonical continuous family** (e.g., connected through the soliton potential \(f\) or through Perelman’s entropy functional)?  
3. In particular, is there any known geometric obstruction (such as anisotropic growth in \(f\)) to uniqueness of the asymptotic cone in this setting?  

Why is it a "good" problem:  
This question captures a deep frontier in the geometric analysis of Ricci shrinkers, focusing on the **uniqueness and analytic rigidity of the asymptotic geometry**—a subtle issue precisely at the intersection of Ricci flow, metric limit theory, and curvature geometry.

- **Profound insight:**  
  The conjecture seeks to understand whether the self‑similar structure of a shrinking soliton determines a unique geometric model at infinity. Uniqueness of the asymptotic cone would amount to an “asymptotic rigidity” principle for solitons, paralleling the uniqueness of tangent flows at singularities in Ricci flow. Its resolution would both deepen our understanding of the analytic stability of self‑similar metrics and clarify whether shrinkers encode a canonical asymptotic geometry.

- **Bridge across fields:**  
  Any progress would require unifying techniques from **Ricci flow analysis** (entropy and monotonicity), **metric geometry** (Gromov–Hausdorff limits and cone structures), and **geometric measure theory** (understanding singular limits and volume ratios). This fusion could open new directions in studying large‑scale geometry of manifolds with lower curvature bounds using PDE tools.

- **Extreme simplicity:**  
  The essential statement —  
  *“Is the asymptotic cone of a noncollapsed nonnegatively curved Ricci shrinker unique?”* —  
  is strikingly concise, geometrically intuitive, yet lies beyond current theory. It encapsulates a central qualitative feature of geometric self‑similarity in an accessible form, illustrating the “narrow entrance, profound depth” ideal of a good problem.